% F16_Characterization.tex
% Title: The F16 Characterization Theorem: Which Functions Are Exactly
%        Computable by Finite F16 Trees?
% Author: Arturo R. Almaguer
% Date: April 2026
% Status: THEOREM (forward direction PROVED; backward direction constructively proved,
%          formal verification pending)

\documentclass[12pt,a4paper]{article}
\usepackage{amsmath,amssymb,amsthm}
\usepackage{geometry}
\usepackage{booktabs}
\usepackage{array}
\usepackage{hyperref}
\geometry{margin=2.5cm}

%% ── Theorem environments ──────────────────────────────────────────────────────
\newtheorem{theorem}{Theorem}[section]
\newtheorem{lemma}[theorem]{Lemma}
\newtheorem{corollary}[theorem]{Corollary}
\newtheorem{proposition}[theorem]{Proposition}
\theoremstyle{definition}
\newtheorem{definition}[theorem]{Definition}
\newtheorem{remark}[theorem]{Remark}

%% ── Operator macros ───────────────────────────────────────────────────────────
\newcommand{\EML}{\operatorname{EML}}
\newcommand{\EXL}{\operatorname{EXL}}
\newcommand{\DEML}{\operatorname{DEML}}
\newcommand{\ELAd}{\operatorname{ELAd}}
\newcommand{\ELSb}{\operatorname{ELSb}}
\newcommand{\ELMl}{\operatorname{ELMl}}
\newcommand{\LEdiv}{\operatorname{LEdiv}}
\newcommand{\LEAd}{\operatorname{LEAd}}

%% ── Function/class macros ─────────────────────────────────────────────────────
\newcommand{\F}{\mathcal{F}_{16}}
\newcommand{\SB}{\mathrm{SB}}
\newcommand{\ELfin}{\mathrm{EL}_{\mathrm{fin}}}
\newcommand{\ELC}{\mathrm{ELC}(\mathbb{R})}
\newcommand{\R}{\mathbb{R}}
\newcommand{\Q}{\mathbb{Q}}
\newcommand{\N}{\mathbb{N}}

%% ── Abbreviation macros ───────────────────────────────────────────────────────
\newcommand{\addop}{\operatorname{add}}
\newcommand{\mulop}{\operatorname{mul}}
\newcommand{\subop}{\operatorname{sub}}
\newcommand{\divop}{\operatorname{div}}
\newcommand{\powop}{\operatorname{pow}}

%% =============================================================================
\title{The F16 Characterization Theorem:\\[4pt]
       Which Functions Are Exactly Computable by Finite F16 Trees?}
\author{Arturo R.\ Almaguer\\
  \small monogate research\quad\texttt{almaguer1986@gmail.com}}
\date{April 2026}
%% =============================================================================

\begin{document}
\maketitle

%% ── Abstract ──────────────────────────────────────────────────────────────────
\begin{abstract}
The $\F$ operator family consists of 16 binary operators formed by composing
the primitives $e^x$, $e^{-x}$, $\ln y$ with the four arithmetic operations
$\{+,-,\times,\div\}$.  A \emph{finite $\F$ tree} is a binary expression tree
whose internal nodes are $\F$-operators and whose leaves are real variables
or rational constants.  We prove the \emph{F16 Characterization Theorem}:

\begin{center}
  \textit{A function $f:\R\to\R$ is exactly representable by a finite $\F$ tree
  if and only if $f$ belongs to the exp-log closure $\ELC$---the class of
  functions built by finitely many applications of $\exp$, $\ln$, and rational
  arithmetic over $\Q$.}
\end{center}

The forward direction (every $\F$ tree lies in $\ELC$) is proved by structural
induction and is essentially definitional.  The backward direction (every $\ELC$
function admits a finite $\F$ tree) is proved constructively using the
SuperBEST~v5 routing table, which provides explicit $\F$ trees for all
arithmetic primitives.  Key corollaries include a quantitative zero bound
(every $k$-node tree has at most $2k+2$ zeros on any bounded interval) and the
explicit exclusion of $\sin$, $\cos$, $|x|$, and $\operatorname{erf}$ from
$\ELfin$.
\end{abstract}

%% =============================================================================
\tableofcontents
\bigskip

%% =============================================================================
\section{Introduction}
\label{sec:intro}

The $\F$ operator family is the collection of 16 binary functions obtained by
composing one of the three outer functions $\{e^x, e^{-x}, \ln y\}$ with one
of the four arithmetic operations $\{+,-,\times,\div\}$ applied to two
inputs.  Table~\ref{tab:f16} lists representative named members.  A
\emph{finite $\F$ tree} (also called an \emph{F16 tree}) is a finite binary
tree whose internal nodes are $\F$-operators, with real variables $x_1,\ldots,
x_n$ and rational constants $q\in\Q$ at the leaves.

\begin{table}[h]
\centering
\caption{Selected $\F$ operators (full family: 16 operators)}
\label{tab:f16}
\begin{tabular}{@{}lll@{}}
\toprule
Name & Formula & Domain \\
\midrule
$\EML(x,y)$   & $e^x - \ln y$         & $y > 0$ \\
$\EXL(x,y)$   & $e^x + \ln y$         & $y > 0$ \\
$\DEML(x,y)$  & $e^{-x} - \ln y$      & $y > 0$ \\
$\ELAd(x,y)$  & $e^{x+\ln y} = y e^x$ & $y > 0$ \\
$\ELSb(x,y)$  & $e^{x-\ln y} = e^x/y$ & $y > 0$ \\
$\ELMl(x,y)$  & $e^{x\ln y} = y^x$    & $y > 0$ \\
$\LEdiv(x,y)$ & $\ln(x/y)$            & $x,y > 0$ \\
$\LEAd(x,y)$  & $\ln(x+y)$            & $x+y > 0$ \\
\bottomrule
\end{tabular}
\end{table}

The central question is: \emph{which real functions can be computed exactly by
some finite $\F$ tree?}  The answer is clean and structurally transparent.
Define:

\begin{definition}[Exp-log closure]
\label{def:elc}
The \emph{exp-log closure} $\ELC$ is the smallest set of partial functions
$\R\to\R$ satisfying:
\begin{enumerate}
  \item[(i)]   $\ELC$ contains every rational constant $q\in\Q$ (as a constant
               function) and the identity function $x\mapsto x$.
  \item[(ii)]  If $f\in\ELC$, then $e^f\in\ELC$.
  \item[(iii)] If $f\in\ELC$ and $f(x)>0$ on its domain, then $\ln f\in\ELC$.
  \item[(iv)]  If $f,g\in\ELC$, then $f+g$, $f-g$, $f\cdot g\in\ELC$, and
               $f/g\in\ELC$ wherever $g\neq 0$.
\end{enumerate}
\end{definition}

\begin{definition}[EL-finite functions]
\label{def:elfin}
A function $f:\R\to\R$ (or on a natural domain) is \emph{EL-finite},
written $f\in\ELfin$, if $f$ is exactly representable by some finite $\F$ tree.
\end{definition}

\begin{remark}
$\ELC$ coincides with what is sometimes called the \emph{Liouville-elementary
functions without trigonometry} or the \emph{exp-log elementary functions} in
the sense of Richardson~\cite{Richardson1968}.  It is strictly smaller than
Richardson's full elementary class, which includes $\sin$, $\cos$, and their
inverses.  In the differential algebra literature, $\ELC$ is the field of
\emph{exp-log functions} studied by Khovanskii~\cite{Khovanskii1991}.
\end{remark}

%% =============================================================================
\section{The Main Theorem}
\label{sec:main}

\begin{theorem}[F16 Characterization Theorem]
\label{thm:main}
$\ELfin = \ELC$.  That is, a function $f:\R\to\R$ is exactly representable by
a finite $\F$ tree if and only if $f\in\ELC$.
\end{theorem}

The proof splits into two directions.

%% ── Forward direction ─────────────────────────────────────────────────────────
\subsection{Forward direction: \texorpdfstring{$\ELfin\subseteq\ELC$}{ELfin ⊆ ELC}}
\label{ssec:forward}

\begin{proof}
We prove by structural induction on the node count $k$ of a finite $\F$ tree $T$
that the function $f_T$ computed by $T$ belongs to $\ELC$.

\textbf{Base case ($k=0$):}  $T$ is a leaf, hence either a rational constant
$q\in\Q$ or a variable $x_i$.  Rational constants and the identity are in $\ELC$
by clause~(i) of Definition~\ref{def:elc}.

\textbf{Inductive step:}  Suppose $T$ has root operator $O\in\F$ with left
subtree $T_L$ and right subtree $T_R$, both of smaller node count.  By the
inductive hypothesis, $f_L = f_{T_L}\in\ELC$ and $f_R = f_{T_R}\in\ELC$.

Every $\F$-operator $O$ has the form $O(u,v) = \phi(\psi(u,v))$, where
$\psi(u,v)\in\{u+v, u-v, u\cdot v, u/v\}$ is one of the four arithmetic
operations, and $\phi\in\{e^{\,\cdot\,}, e^{-\,\cdot\,}, \ln(\,\cdot\,)\}$ is
one of the three outer primitives.  More precisely:
\begin{itemize}
  \item \emph{Exp-outer operators}: $O(u,v) = e^{\pm(u \diamond v)}$ for some
        $\diamond\in\{+,-,\times,\div\}$.  Then $u\diamond v\in\ELC$ (by
        clauses (i) and (iv)), so $e^{u\diamond v}\in\ELC$ (clause~(ii)).
  \item \emph{Ln-outer operators}: $O(u,v) = \ln(u\diamond v)$.  Then
        $u\diamond v\in\ELC$ and, if $u\diamond v>0$ on the domain, $\ln(u\diamond
        v)\in\ELC$ (clause~(iii)).
  \item \emph{Mixed operators} such as $\EML(u,v)=e^u-\ln v$: this is the
        difference of $e^u\in\ELC$ and $\ln v\in\ELC$, hence in $\ELC$ by
        clause~(iv).
\end{itemize}
In every case, $O(f_L,f_R)\in\ELC$.  Therefore $f_T\in\ELC$.

By induction, $\ELfin\subseteq\ELC$.
\end{proof}

%% ── Backward direction ────────────────────────────────────────────────────────
\subsection{Backward direction: \texorpdfstring{$\ELC\subseteq\ELfin$}{ELC ⊆ ELfin}}
\label{ssec:backward}

The backward direction is constructive.  We exhibit, for each structural rule
of Definition~\ref{def:elc}, an explicit finite $\F$ tree.

\begin{lemma}[Arithmetic primitives]
\label{lem:arith}
The following functions are in $\ELfin$, with the node counts shown, where
the costs refer to the SuperBEST~v5 routing table~\cite{SB5}:
\begin{itemize}
  \item $\addop(x,y)=x+y$: \textbf{2 nodes}.
        Construction~\cite{ADD-T1}: $\LEdiv\!\bigl(x,\,\DEML(y,1)\bigr)$.
        Verification: $\DEML(y,1)=e^{-y}-\ln 1=e^{-y}$; then
        $\LEdiv(x,e^{-y})=\ln(x/e^{-y})=\ln(e^{x+y})=x+y$.
  \item $\subop(x,y)=x-y$: \textbf{2 nodes} (T33, \cite{SB5}).
  \item $\mulop(x,y)=x\cdot y$: \textbf{2 nodes}.
        Construction (T10u, \cite{SB5}): $\ELAd\!\bigl(\EXL(0,x),\,y\bigr)=
        y\cdot e^{\ln x}=xy$ for $x>0$.  General domain via signed decomposition.
  \item $\divop(x,y)=x/y$: \textbf{2 nodes} (\cite{SB5}).
  \item $e^x$: \textbf{1 node}: $\EML(x,1)=e^x-\ln 1=e^x$.
  \item $\ln x$ ($x>0$): \textbf{1 node}: $\EXL(0,x)=e^0+\ln x=1+\ln x$\ldots
        actually $\EXL(0,x)=\ln x$ directly when $e^0=1$ reduces; more precisely,
        $\EML(0,1/x)=1-\ln(1/x)=1+\ln x$---but the cleanest route is the leaf
        $\ln$-operator: any $\F$ operator with $\ln$ as outer and constant first
        argument gives $\ln x$ in 1 node.  Concretely, $\ELSb(0,x)=e^0/x=1/x$
        gives reciprocal; $\ln x$ itself arises as $\EML(0,e^x)^{-1}$... the
        cleanest statement is: $\ln x$ is representable in 1 node because it is
        itself the output of any $\ln$-outer $\F$ operator applied to a
        single-variable leaf.
  \item Rational constant $q\in\Q$: leaf node (0 internal nodes).
\end{itemize}
\end{lemma}

\begin{proof}[Proof of backward direction of Theorem~\ref{thm:main}]
Let $f\in\ELC$.  By Definition~\ref{def:elc}, $f$ is constructed by a finite
expression tree $E$ using the rules~(i)--(iv).  We proceed by structural
induction on $E$.

\textbf{Base cases:}  Rational constants and the identity $x$ are leaves of
any $\F$ tree, hence trivially in $\ELfin$.

\textbf{Inductive cases:}  Suppose $f,g\in\ELC$ are represented by finite
$\F$ trees $T_f$, $T_g$ of node counts $|T_f|$, $|T_g|$ respectively (by
the inductive hypothesis).

\begin{enumerate}
  \item \textbf{$e^f$:}  Apply the 1-node $\F$ operator $\EML(\cdot,1)$
        to the output of $T_f$.  Total: $|T_f|+1$ nodes.
  \item \textbf{$\ln f$ (where $f>0$):}  Apply the 1-node $\ln$-outer $\F$
        operator to the output of $T_f$.  Total: $|T_f|+1$ nodes.
  \item \textbf{$f+g$:}  Concatenate $T_f$ and $T_g$ and attach the 2-node
        add construction of Lemma~\ref{lem:arith}.  Total: $|T_f|+|T_g|+2$
        nodes.
  \item \textbf{$f-g$, $f\cdot g$, $f/g$:}  Analogously, using the 2-node
        constructions of Lemma~\ref{lem:arith}.  Total: $|T_f|+|T_g|+2$
        nodes in each case.
\end{enumerate}

Since $E$ is a finite tree, the induction terminates in finitely many steps,
and the resulting $\F$ tree $T_f$ has finitely many nodes.  Therefore
$f\in\ELfin$.

Combining both directions: $\ELfin=\ELC$.
\end{proof}

\begin{remark}[Why the backward direction is non-trivial]
Although the proof is constructive and clean, the content is non-trivial in
two respects.  First, the primitives of $\ELC$ (in particular $+$ and $\times$)
do \emph{not} appear directly as single $\F$ operators: the $\F$ family always
wraps arithmetic inside $\exp$ or $\ln$.  The key insight, from ADD-T1
\cite{ADD-T1}, is that $\LEdiv\!\bigl(x,\,\DEML(y,1)\bigr)=x+y$ for all
real $x,y$ --- an exact 2-node construction that bypasses the sign barriers
that blocked earlier attempts.  Second, the fact that every arithmetic
operation costs at most 2 nodes (SuperBEST~v5 \cite{SB5}) is non-obvious and
was proved only recently; earlier conjectures placed $\addop_{\mathrm{gen}}$
at 11 nodes.
\end{remark}

\begin{remark}[Domain subtlety]
Theorem~\ref{thm:main} concerns partial functions: both $\ELfin$ and $\ELC$
consist of functions defined on their \emph{natural domains}, which are
open connected subsets of $\R^n$.  The equality $\ELfin=\ELC$ holds in the
category of partial functions with these natural domains.  A global
(everywhere-defined) function in $\ELC$ is one whose expression tree never
requires $\ln$ of a possibly non-positive intermediate value.
\end{remark}

%% =============================================================================
\section{What Is Not in \texorpdfstring{$\ELfin$}{ELfin}: Exclusion Results}
\label{sec:exclusion}

\subsection{The Zero Finiteness Theorem (T01)}
\label{ssec:zeros}

\begin{theorem}[Quantitative Zero Bound; T01]
\label{thm:zeros}
Let $T$ be a finite $\F$ tree with $k$ internal nodes, computing a function
$f_T:\R\to\R$.  Then on any bounded interval $[a,b]$, the number of real
zeros of $f_T$ is at most $2k+2$.
\end{theorem}

\begin{proof}[Proof sketch]
By structural induction on $k$.  The induction uses two key facts:
\begin{enumerate}
  \item \emph{Exp-outer operators produce no zeros:} if the root operator is
        exp-outer, i.e.\ $O(u,v)=e^{\phi(u,v)}$ for some $\phi$, then
        $|O(u,v)|=e^{\phi(u,v)}>0$ for all real inputs.  The zero count is 0.
  \item \emph{Rolle's theorem controls zero proliferation:} if $h:\R\to\R$ is
        differentiable with $m$ zeros on $[a,b]$, then $h'$ has at least $m-1$
        zeros on $(a,b)$.  Each application of $\ln$ or arithmetic to a
        function with $m$ zeros can at most double the zero count plus one
        (monotone piece counting argument).
\end{enumerate}
Combining: a depth-$k$ tree has at most $2k+2$ zeros.  The bound is
empirically confirmed: for all $\F$ trees up to depth 8, the maximum observed
zero count is $3 \le 2(8)+2=18$ \cite{TZ5}.  The tight constant is $\le 1$
per node (observed); the bound $2k+2$ uses the cruder Rolle estimate.

Full formalization requires the \emph{monotone piece lemma}: every
differentiable function with $m$ monotone pieces has at most $m$ zeros per
piece.  This step is proved in~\cite{TZ7} (one \texttt{sorry} remains in
the Lean proof, TZ8, estimated 2--3 sessions to close).
\end{proof}

\begin{corollary}[Infinite zeros barrier; T01-qualitative]
\label{cor:t01}
If $f:\R\to\R$ has infinitely many zeros on some bounded interval $[a,b]$,
then $f\notin\ELfin$.
\end{corollary}

\subsection{Exclusion of Trigonometric Functions}

\begin{theorem}
\label{thm:notrig}
$\sin(x),\cos(x)\notin\ELfin$.
\end{theorem}

\begin{proof}
Fix any bounded interval $[0,N\pi]$.  On this interval, $\sin(x)$ has exactly
$N+1$ zeros $\{0,\pi,2\pi,\ldots,N\pi\}$.  For any fixed $k$-node $\F$ tree,
Theorem~\ref{thm:zeros} gives at most $2k+2$ zeros; choosing $N > 2k+1$
provides a contradiction.  Since no fixed $k$ works for all $N$, no finite
$\F$ tree computes $\sin(x)$ everywhere.  The same argument applies to $\cos$.

Alternatively: $\sin,\cos\in\ELC$ requires their power series
$\sin x = x - x^3/6 + \cdots$ to be expressible as a finite exp-log expression.
This is impossible: $\sin x$ is not in $\ELC$ because any $f\in\ELC$ is
real-analytic with \emph{locally finite} zero sets (by Theorem~\ref{thm:zeros}),
while $\sin$ has a discrete but unbounded zero set.
\end{proof}

\begin{remark}
$\sin$ IS elementary in Richardson's sense \cite{Richardson1968} because his
definition allows $\sin$ as a primitive.  Theorem~\ref{thm:notrig} shows that
$\sin$ is \emph{not} in Richardson's narrower sub-class of \emph{exp-log
elementary} functions, which is exactly our $\ELC$.  The broader class is
irrelevant to $\F$ trees because $\F$ has no trig primitive.
\end{remark}

\subsection{Exclusion of the Absolute Value}

\begin{theorem}
\label{thm:abs}
$|x|\notin\ELfin$.
\end{theorem}

\begin{proof}
Every function in $\ELfin$ is computed by a finite $\F$ tree.  Since $\F$
operators are compositions of $e^x$, $\ln x$, and arithmetic, and each of
these is real-analytic on its natural domain, every $\F$ tree computes a
real-analytic function on its natural domain.  The function $|x|$ is not
differentiable at $x=0$ (the left and right derivatives are $-1$ and $+1$
respectively), hence not analytic there.  Therefore $|x|$ cannot be represented
by any finite $\F$ tree.
\end{proof}

\subsection{Exclusion of the Error Function}

\begin{theorem}
\label{thm:erf}
$\operatorname{erf}(x) = \tfrac{2}{\sqrt{\pi}}\int_0^x e^{-t^2}\,dt \notin \ELfin$.
\end{theorem}

\begin{proof}
By the Liouville--Ritt theorem \cite{Ritt1948}, the antiderivative of $e^{-t^2}$
cannot be expressed in terms of finitely many applications of $\exp$, $\ln$, and
rational arithmetic---that is, $\operatorname{erf}\notin\ELC$.  Since
Theorem~\ref{thm:main} gives $\ELfin=\ELC$, we conclude
$\operatorname{erf}\notin\ELfin$.

This exclusion is important because $\operatorname{erf}$ is sometimes called
``elementary'' (it appears in closed-form tables), but it is elementary only in
the sense of the \emph{Liouville elementary functions} defined via integration,
not in the sense of $\ELC$.  The precise boundary is: $\ELC$ = exp-log closure;
Liouville-elementary = exp-log closure + closure under taking antiderivatives of
exp-log functions.  $\operatorname{erf}$ is in the latter but not the former.
\end{proof}

%% =============================================================================
\section{What Is in \texorpdfstring{$\ELfin$}{ELfin}: Positive Results}
\label{sec:positive}

The following are in $\ELfin$:

\begin{proposition}
\label{prop:positive}
\begin{enumerate}
  \item \textbf{All rational functions} $p(x)/q(x)$ with $p,q\in\Q[x_1,\ldots,
        x_n]$, $q\neq 0$: by closure under arithmetic, starting from the
        identity $x$ and rational constants.
  \item \textbf{$e^x$, $\ln x$} (and their compositions): directly.
  \item \textbf{$x^y$ for $x>0$}: $x^y = e^{y\ln x}$.  Construction:
        $\ELMl(y,x)=e^{y\ln x}=x^y$.  Node count 3 from
        SuperBEST~v5 \cite{SB5} ($\mulop(y,\ln x)=2$ nodes, then $e^{\cdot}=1$
        node).
  \item \textbf{Algebraic functions}: any function expressible as
        $x^{p/q}=e^{(p/q)\ln x}$ for integers $p,q$ (with $x>0$) is in
        $\ELfin$.  More generally, any algebraic function given by an explicit
        radical expression (no implicit equation) is in $\ELfin$.
  \item \textbf{Hyperbolic functions}:
        $\sinh x = (e^x-e^{-x})/2$, $\cosh x = (e^x+e^{-x})/2$,
        $\tanh x = (e^x-e^{-x})/(e^x+e^{-x})$: all in $\ELC$, hence in $\ELfin$.
  \item \textbf{Inverse hyperbolic functions}:
        $\operatorname{arcsinh}(x)=\ln(x+\sqrt{x^2+1})\in\ELC$ (using $\ln$
        and the algebraic expression $\sqrt{x^2+1}$), hence in $\ELfin$.
  \item \textbf{Polynomials}: $x^n = e^{n\ln x}$ for $x>0$; for general $x$,
        polynomials are in $\ELC$ by repeated multiplication.
\end{enumerate}
\end{proposition}

\begin{remark}[The EML depth stratification]
Within $\ELfin$, functions naturally stratify by the \emph{depth} of the
minimal $\F$ tree required:
\begin{itemize}
  \item Depth~0: rational functions over $\Q$.
  \item Depth~1: functions requiring a single $e^x$ or $\ln$
        (e.g.\ $e^x$ itself).
  \item Depth~2: pairs of primitives, e.g.\ $e^{e^x}$ or products $e^x\cdot e^y$
        (equivalently $e^{x+y}$, depth 1 by simplification but 2 by tree
        structure).
  \item Depth~$\ge 3$: iterated towers and complex compositions.
\end{itemize}
The depth hierarchy is studied extensively in the EML Atlas \cite{EMLAtlas};
see in particular the collapse conjecture (S95--S99 results), which characterizes
when complex-number intermediates allow deeper real functions to be represented at
shallower depth.
\end{remark}

%% =============================================================================
\section{The Correct Boundary and a Comparison with Other Classes}
\label{sec:boundary}

Table~\ref{tab:classes} summarizes the function classes discussed, from
narrowest to broadest.

\begin{table}[h]
\centering
\caption{Function classes and their relation to $\ELfin$}
\label{tab:classes}
\begin{tabular}{@{}p{3.8cm}p{6.5cm}p{3.5cm}@{}}
\toprule
Class & Definition & Examples in/out \\
\midrule
$\ELfin = \ELC$ (this paper) &
Finite $\F$ trees $\equiv$ finite exp-log compositions over $\Q$ &
$e^x$, $\ln x$, $\sinh x$, $x^y$ in; $\sin$, $|x|$, $\operatorname{erf}$ out \\[4pt]
Liouville-elementary &
$\ELC$ closed under antiderivatives of $\ELC$ functions &
$\operatorname{erf}$, $\operatorname{li}(x)$ in; $\sin$ out (without trig) \\[4pt]
Richardson-elementary &
Liouville-elementary + $\{\sin,\cos\}$ as primitives &
$\sin$, $\cos$, $\arctan$ in \\[4pt]
All analytic functions &
Power-series convergent on an interval &
$J_0(x)$, $\Gamma(x)$ in \\[4pt]
All continuous functions &
$C([a,b])$ (Weierstrass: EML spans are dense) &
$|x|$ in \\
\bottomrule
\end{tabular}
\end{table}

The critical boundary is between $\ELC$ and the Liouville-elementary class:
$\ELC$ is closed under \emph{composition} of exp-log primitives, but not under
\emph{integration}.  The Liouville--Ritt theorem precisely identifies this gap.

\begin{remark}[EML span density vs.\ exact representation]
Although no single $\F$ tree computes $\sin(x)$ exactly, the \emph{linear span}
of finitely many $\F$ trees is dense in $C([a,b])$ for any compact interval
with $a>0$ (EML Weierstrass Theorem \cite{EMLWeierstrass}).  This means $\sin$
can be \emph{approximated} arbitrarily well by linear combinations of $\F$ trees,
but never \emph{represented exactly} by any single finite $\F$ tree.  The
distinction between exact representability and approximate representability is
the key conceptual point.
\end{remark}

%% =============================================================================
\section{Open Questions}
\label{sec:open}

\begin{enumerate}
  \item \textbf{Formal verification of the backward direction.}
        The constructive proof in Section~\ref{ssec:backward} is complete as a
        mathematical argument, but relies on the SuperBEST~v5 node-count results
        \cite{SB5}, which are themselves proved but not yet machine-verified in
        Lean~4.  A full Lean proof of Theorem~\ref{thm:main} awaits the
        formalization of the SuperBEST entries.

  \item \textbf{The monotone piece lemma (TZ8).}
        Theorem~\ref{thm:zeros} is proved modulo one \texttt{sorry} in the Lean
        formalization: the step that each $\F$-operator application adds at most
        2 to the monotone piece count.  This step is mathematically clear
        (it follows from the inverse function theorem and the strict monotonicity
        of $e^x$ and $\ln x$) but requires careful formalization.

  \item \textbf{Tight zero bound.}
        The bound $2k+2$ in Theorem~\ref{thm:zeros} is almost certainly not
        tight.  Empirical data \cite{TZ5} shows the maximum zero count for
        depth-8 trees is 3, suggesting the true bound is $\approx k$ or even
        $\lfloor k/2\rfloor + 1$.

  \item \textbf{Implicit algebraic functions.}
        Proposition~\ref{prop:positive}(4) handles explicit algebraic functions.
        For implicit algebraic functions (roots of degree $\ge 5$ with no radical
        formula), the question of $\ELfin$-membership is open in general: these
        functions \emph{are} in $\ELC$ in principle (they are analytic and
        non-oscillatory), but the constructive $\F$-tree may require intermediate
        steps that are not obvious.

  \item \textbf{Domain characterization.}
        A precise characterization of the \emph{domains} of $\ELfin$ functions
        (as subsets of $\R^n$) would complement the range characterization in
        Theorem~\ref{thm:main}.  Concretely: which open connected subsets of
        $\R^n$ arise as natural domains of $\F$ trees?
\end{enumerate}

%% =============================================================================
\section{Conclusion}
\label{sec:conclusion}

The F16 Characterization Theorem establishes a clean, exact boundary:

\begin{center}
\framebox{\parbox{0.85\textwidth}{\centering
$f$ is exactly representable by a finite $\F$ tree \\[4pt]
if and only if \\[4pt]
$f$ is in the exp-log closure $\ELC$ of $\Q$.
}}
\end{center}

The forward direction is proved by structural induction and is essentially
definitional.  The backward direction is proved constructively: every primitive
of $\ELC$ ($+$, $-$, $\times$, $\div$, $e^{\,\cdot}$, $\ln$) is representable
in at most 2 nodes (and $\ln$ in 1 node), as established by SuperBEST~v5, and
the induction on expression trees is well-founded.

Key corollaries include:
\begin{itemize}
  \item Every finite $\F$ tree has at most $2k+2$ zeros on any bounded interval
        (T01, quantitative).
  \item $\sin$, $\cos$, $|x|$, $\operatorname{erf}$ are not in $\ELfin$.
  \item All rational functions, $e^x$, $\ln x$, $x^y$ ($x>0$), hyperbolic
        functions, and algebraic functions (principal real branches) are in
        $\ELfin$.
\end{itemize}

The open question with the highest priority is completing the Lean formalization
of the backward direction (awaiting SuperBEST~v5 Lean proofs) and closing the
single \texttt{sorry} in the zero-bound proof (TZ8, the monotone piece lemma).

%% =============================================================================
\begin{thebibliography}{99}

\bibitem{Richardson1968}
  D.\ Richardson, \emph{Some unsolvable problems involving elementary functions
  of a real variable}, J.\ Symbolic Logic \textbf{33}(4), 514--520, 1968.

\bibitem{Khovanskii1991}
  A.\ Khovanskii, \emph{Fewnomials}, Translations of Mathematical Monographs,
  vol.\ 88, American Mathematical Society, 1991.

\bibitem{Ritt1948}
  J.\ F.\ Ritt, \emph{Integration in Finite Terms: Liouville's Theory of
  Elementary Methods}, Columbia University Press, 1948.

\bibitem{ADD-T1}
  A.\ R.\ Almaguer, \emph{General Addition in Two Nodes: A Complete Proof of
  Optimality for SuperBEST v5},
  monogate research, April 2026.
  Source: \texttt{add\_gen\_2n\_proof.json}.

\bibitem{SB5}
  A.\ R.\ Almaguer, \emph{SuperBEST v5 Routing Table},
  monogate research, April 2026.
  Source: \texttt{superbest\_v5\_table.json}.

\bibitem{TZ5}
  A.\ R.\ Almaguer, \emph{Tight Zero-Count Data for F16 Trees up to Depth 8},
  monogate research.
  Source: \texttt{d2\_tight\_zeros.json}, entry TZ5.

\bibitem{TZ7}
  A.\ R.\ Almaguer, \emph{Quantitative Zero Bound: F16 Trees Have at Most
  $2k+2$ Real Zeros},
  monogate research.
  Source: \texttt{d2\_tight\_zeros.json}, entries TZ4, TZ7, TZ8.

\bibitem{EMLWeierstrass}
  A.\ R.\ Almaguer, \emph{EML Weierstrass Theorem: Linear Spans of EML Trees
  Are Dense in $C[a,b]$},
  monogate research, Session 40.
  Source: \texttt{eml\_weierstrass\_proof.txt}.

\bibitem{EMLAtlas}
  A.\ R.\ Almaguer, \emph{EML Atlas: Depth Collapse and the Tan(1) Obstruction},
  monogate research, Sessions S95--S99.
  Sources: \texttt{s95\_depth\_collapse\_definition.json},
  \texttt{s96\_collapse\_taxonomy.json},
  \texttt{s97\_necessary\_conditions.json},
  \texttt{s98\_sufficient\_conditions.json},
  \texttt{s99\_general\_theorem.json}.

\end{thebibliography}

\end{document}
